\documentclass[11pt,a4paper]{article}
\usepackage[utf8]{inputenc}
\usepackage[T1]{fontenc}
\usepackage[catalan]{babel}
\usepackage{lmodern}
\usepackage[sfdefault,lining]{FiraSans}
\usepackage[a4paper,top=2.3cm,bottom=2.2cm,left=1.8cm,right=1.8cm,headheight=1.3cm,headsep=0.4cm,footskip=1cm]{geometry}
\usepackage{amsmath,amssymb}
\allowdisplaybreaks
\usepackage{icomma}
\usepackage{graphicx}
\usepackage[table]{xcolor}
\usepackage{tikz}
\usepackage{fontawesome5}
\usepackage{ccicons}
\usepackage{fancyhdr}
\usepackage{lastpage}
\usepackage{tasks}
\usepackage{needspace}
\usepackage[most]{tcolorbox}
\usepackage{array}
\usepackage[hidelinks,pdftitle={Fitxa d'exercicis de la unitat 3: Vectors i aplicacions. Rectes, plans i feixos de plans. 2n de Batxillerat MII},pdfsubject={Matemàtiques II},pdfkeywords={vectors, producte escalar, producte vectorial, producte mixt, àrees, volums, equacions de la recta, equacions del pla, posicions relatives, feixos de plans, PAU}]{hyperref}

\definecolor{navy}{HTML}{1B2748}
\definecolor{accent}{HTML}{E8603A}
\definecolor{teal}{HTML}{1F9E90}
\definecolor{soft}{HTML}{EEF1F8}
\definecolor{muted}{HTML}{5D6685}

\setlength{\parindent}{0pt}
\setlength{\parskip}{3pt}
\settasks{label=\alph*), label-format=\bfseries\color{navy}, label-width=1.3em, item-indent=1.8em, column-sep=0.6em, after-item-skip=5pt, before-skip=3pt, after-skip=4pt}

% ---------- capçalera i peu ----------
\newcommand{\llicencia}{\href{https://creativecommons.org/licenses/by-nc-sa/4.0/deed.ca}{\ccbyncsa}\ \ Llicència Creative Commons Reconeixement-NoComercial-CompartirIgual 4.0 Internacional (CC BY-NC-SA 4.0)}
\pagestyle{fancy}
\fancyhf{}
\fancyhead[L]{\raisebox{-0.35\height}{\includegraphics[height=0.95cm]{logo-gva.png}}}
\fancyhead[C]{\small\color{muted}Fitxa d'exercicis · Unitat 3: Vectors, rectes i plans · 2n Batxillerat MII}
\fancyhead[R]{\raisebox{-0.35\height}{\includegraphics[height=0.95cm]{logo-vermellar.png}}}
\fancyfoot[L]{\scriptsize\color{muted}\llicencia}
\fancyfoot[R]{\small\color{muted}\thepage}
\renewcommand{\headrulewidth}{0.4pt}
\renewcommand{\headrule}{\color{muted!40}\hrule height \headrulewidth}
\fancypagestyle{primera}{\fancyhf{}\renewcommand{\headrulewidth}{0pt}
  \fancyfoot[L]{\scriptsize\color{muted}\llicencia}
  \fancyfoot[R]{\small\color{muted}\thepage}}

% ---------- nivells i exercicis ----------
\newcommand{\estrelles}[1]{{\color{accent}\ifcase#1\or\faStar\faStar[regular]\faStar[regular]\or\faStar\faStar\faStar[regular]\or\faStar\faStar\faStar\fi}}
\newcounter{exer}
\newcommand{\ex}[1]{\par\needspace{4\baselineskip}\medskip\refstepcounter{exer}\noindent{\color{navy}\bfseries\theexer.}\ \ }

\newcommand{\seccio}[2]{\par\needspace{10\baselineskip}\bigskip
  \begin{tcolorbox}[colback=navy,colframe=navy,coltext=white,arc=2mm,boxrule=0pt,left=3mm,right=3mm,top=1.2mm,bottom=1.2mm]
  {\large\bfseries #1\quad #2}
  \end{tcolorbox}\vspace{-2pt}}


\setlength{\emergencystretch}{3em}  % les fórmules llargues (equacions de rectes, ternes) passen a la línia següent

% quadre de recordatori (només a la fitxa d'exercicis)
\newcommand{\recorda}[1]{\par\smallskip\begin{tcolorbox}[colback=accent!7,colframe=accent!55,arc=2mm,boxrule=0.5pt,left=2.5mm,right=2.5mm,top=1mm,bottom=1mm]\small{\color{accent}\faLightbulb[regular]}\ \textbf{Recorda:}\ #1\end{tcolorbox}}
% etiqueta dels exercicis de la PAU d'altres comunitats
\newcommand{\pau}[1]{{\color{teal}\small\bfseries\faUniversity\ PAU · #1.}\ }
\newcommand{\res}[1]{#1}\newcommand{\rest}[1]{#1}

% =====================================================================
% Capçalera FentMates de les fitxes: un color per a cada curs.
% Els colors són els de la pàgina inicial (LEVELS a assets/js/materials.js);
% el logo ve de _eines/logo/pdf/ (python3 _eines/logo/genera.py).
%
% Ús, al final del preàmbul de la unitat:
%   % =====================================================================
% Capçalera FentMates de les fitxes: un color per a cada curs.
% Els colors són els de la pàgina inicial (LEVELS a assets/js/materials.js);
% el logo ve de _eines/logo/pdf/ (python3 _eines/logo/genera.py).
%
% Ús, al final del preàmbul de la unitat:
%   % =====================================================================
% Capçalera FentMates de les fitxes: un color per a cada curs.
% Els colors són els de la pàgina inicial (LEVELS a assets/js/materials.js);
% el logo ve de _eines/logo/pdf/ (python3 _eines/logo/genera.py).
%
% Ús, al final del preàmbul de la unitat:
%   \input{../comu/fentmates.tex}
%   \fmcurs{1eso}
% i, al començament del document:
%   \capcalera{Matemàtiques}{1r ESO}{Fitxa d'exercicis · Unitat 2}{Divisibilitat}
% Opcional abans de \input: \def\fmneed{14} (espai mínim abans de cada secció).
% =====================================================================
\makeatletter
\def\fm@logodir{../../logo/pdf}
\providecommand{\fmneed}{10}
\def\fm@cursnom{}

% clau del curs, color i color del text damunt del color
\newcommand{\fm@curs}[3]{\@namedef{fm@color@#1}{#2}\@namedef{fm@text@#1}{#3}}
\fm@curs{1eso}{E6B84A}{navy}
\fm@curs{2eso}{3B82F6}{white}
\fm@curs{3eso}{FF7A50}{white}
\fm@curs{4eso-a}{A78BFA}{white}
\fm@curs{4eso-b}{6D28D9}{white}
\fm@curs{1batx-mi}{9B1C1C}{white}
\fm@curs{1batx-ccss}{EF6B6B}{white}
\fm@curs{2batx-mii}{1E6B47}{white}
\fm@curs{2batx-ccss}{2BB3A3}{white}

\newcommand{\fmcurs}[1]{%
  \@ifundefined{fm@color@#1}{\PackageError{fentmates}{Curs desconegut: #1}{}}{}%
  \edef\fm@tmp{\noexpand\definecolor{curs}{HTML}{\@nameuse{fm@color@#1}}%
               \noexpand\colorlet{curstext}{\@nameuse{fm@text@#1}}}\fm@tmp
  \colorlet{cursfosc}{curs!72!black}%
  \def\fm@logo{\fm@logodir/vertical-#1.pdf}%
  \def\fm@marca{\fm@logodir/marca-#1.pdf}}

% ---------- capçalera de la primera pàgina ----------
% #1 matèria (petit), #2 curs, #3 tipus de document i unitat, #4 títol
\newcommand{\capcalera}[4]{%
  \gdef\fm@cursnom{#2}%
  \par\noindent
  \begin{tcolorbox}[colback=curs,colframe=curs,boxrule=0pt,arc=2.2mm,width=3.5cm,height=2.15cm,
      valign=center,left=3mm,right=2mm,top=1mm,bottom=1mm,nobeforeafter,box align=center]
    \color{curstext}{\footnotesize #1}\par\vspace{1pt}{\fontsize{19}{21}\selectfont\bfseries #2\par}%
  \end{tcolorbox}\hspace{5mm}%
  \begin{minipage}[c]{\dimexpr\linewidth-3.5cm-5mm-2.9cm\relax}
    {\color{cursfosc}\small\bfseries\MakeUppercase{#3}}\par\vspace{2pt}
    {\color{navy}\fontsize{20}{23}\selectfont\bfseries #4\par}
  \end{minipage}\hfill
  \begin{minipage}[c]{2.6cm}\raggedleft\includegraphics[width=2.6cm]{\fm@logo}\end{minipage}\par
  \vspace{2pt}}

% ---------- barres de secció del color del curs ----------
\renewcommand{\seccio}[2]{\par\needspace{\fmneed\baselineskip}\bigskip
  \begin{tcolorbox}[colback=curs,colframe=curs,coltext=curstext,arc=2mm,boxrule=0pt,left=3mm,right=3mm,top=1.2mm,bottom=1.2mm]
  {\large\bfseries #1\quad #2}
  \end{tcolorbox}\vspace{-2pt}}

% ---------- estrelles de dificultat del color del curs (a Batxillerat no se'n posen) ----------
\@ifundefined{estrelles}{}{%
  \renewcommand{\estrelles}[1]{{\color{curs}\ifcase#1\or\faStar\faStar[regular]\faStar[regular]\or\faStar\faStar\faStar[regular]\or\faStar\faStar\faStar\fi}}}

% ---------- capçalera i peu de les pàgines ----------
% primera pàgina: logos institucionals al peu
\fancypagestyle{primera}{\fancyhf{}\renewcommand{\headrulewidth}{0pt}
  \fancyfoot[L]{\raisebox{-0.32\height}{\includegraphics[height=0.78cm]{logo-gva.png}}\hspace{5mm}%
    \raisebox{-0.32\height}{\includegraphics[height=0.78cm]{logo-vermellar.png}}}
  \fancyfoot[R]{\scriptsize\color{muted}\href{https://creativecommons.org/licenses/by-nc-sa/4.0/deed.ca}{\ccbyncsa}\ \ CC BY-NC-SA 4.0 · fentmates.org\qquad\small\thepage}}
% la resta: el símbol a l'esquerra i el curs a la dreta, amb el títol al centre com abans
\fancyhead[L]{\raisebox{-0.3\height}{\includegraphics[height=0.62cm]{\fm@marca}}}
\fancyhead[R]{\small\bfseries\color{cursfosc}\fm@cursnom}
\renewcommand{\headrule}{\color{curs!45}\hrule height \headrulewidth}
\makeatother

%   \fmcurs{1eso}
% i, al començament del document:
%   \capcalera{Matemàtiques}{1r ESO}{Fitxa d'exercicis · Unitat 2}{Divisibilitat}
% Opcional abans de \input: \def\fmneed{14} (espai mínim abans de cada secció).
% =====================================================================
\makeatletter
\def\fm@logodir{../../logo/pdf}
\providecommand{\fmneed}{10}
\def\fm@cursnom{}

% clau del curs, color i color del text damunt del color
\newcommand{\fm@curs}[3]{\@namedef{fm@color@#1}{#2}\@namedef{fm@text@#1}{#3}}
\fm@curs{1eso}{E6B84A}{navy}
\fm@curs{2eso}{3B82F6}{white}
\fm@curs{3eso}{FF7A50}{white}
\fm@curs{4eso-a}{A78BFA}{white}
\fm@curs{4eso-b}{6D28D9}{white}
\fm@curs{1batx-mi}{9B1C1C}{white}
\fm@curs{1batx-ccss}{EF6B6B}{white}
\fm@curs{2batx-mii}{1E6B47}{white}
\fm@curs{2batx-ccss}{2BB3A3}{white}

\newcommand{\fmcurs}[1]{%
  \@ifundefined{fm@color@#1}{\PackageError{fentmates}{Curs desconegut: #1}{}}{}%
  \edef\fm@tmp{\noexpand\definecolor{curs}{HTML}{\@nameuse{fm@color@#1}}%
               \noexpand\colorlet{curstext}{\@nameuse{fm@text@#1}}}\fm@tmp
  \colorlet{cursfosc}{curs!72!black}%
  \def\fm@logo{\fm@logodir/vertical-#1.pdf}%
  \def\fm@marca{\fm@logodir/marca-#1.pdf}}

% ---------- capçalera de la primera pàgina ----------
% #1 matèria (petit), #2 curs, #3 tipus de document i unitat, #4 títol
\newcommand{\capcalera}[4]{%
  \gdef\fm@cursnom{#2}%
  \par\noindent
  \begin{tcolorbox}[colback=curs,colframe=curs,boxrule=0pt,arc=2.2mm,width=3.5cm,height=2.15cm,
      valign=center,left=3mm,right=2mm,top=1mm,bottom=1mm,nobeforeafter,box align=center]
    \color{curstext}{\footnotesize #1}\par\vspace{1pt}{\fontsize{19}{21}\selectfont\bfseries #2\par}%
  \end{tcolorbox}\hspace{5mm}%
  \begin{minipage}[c]{\dimexpr\linewidth-3.5cm-5mm-2.9cm\relax}
    {\color{cursfosc}\small\bfseries\MakeUppercase{#3}}\par\vspace{2pt}
    {\color{navy}\fontsize{20}{23}\selectfont\bfseries #4\par}
  \end{minipage}\hfill
  \begin{minipage}[c]{2.6cm}\raggedleft\includegraphics[width=2.6cm]{\fm@logo}\end{minipage}\par
  \vspace{2pt}}

% ---------- barres de secció del color del curs ----------
\renewcommand{\seccio}[2]{\par\needspace{\fmneed\baselineskip}\bigskip
  \begin{tcolorbox}[colback=curs,colframe=curs,coltext=curstext,arc=2mm,boxrule=0pt,left=3mm,right=3mm,top=1.2mm,bottom=1.2mm]
  {\large\bfseries #1\quad #2}
  \end{tcolorbox}\vspace{-2pt}}

% ---------- estrelles de dificultat del color del curs (a Batxillerat no se'n posen) ----------
\@ifundefined{estrelles}{}{%
  \renewcommand{\estrelles}[1]{{\color{curs}\ifcase#1\or\faStar\faStar[regular]\faStar[regular]\or\faStar\faStar\faStar[regular]\or\faStar\faStar\faStar\fi}}}

% ---------- capçalera i peu de les pàgines ----------
% primera pàgina: logos institucionals al peu
\fancypagestyle{primera}{\fancyhf{}\renewcommand{\headrulewidth}{0pt}
  \fancyfoot[L]{\raisebox{-0.32\height}{\includegraphics[height=0.78cm]{logo-gva.png}}\hspace{5mm}%
    \raisebox{-0.32\height}{\includegraphics[height=0.78cm]{logo-vermellar.png}}}
  \fancyfoot[R]{\scriptsize\color{muted}\href{https://creativecommons.org/licenses/by-nc-sa/4.0/deed.ca}{\ccbyncsa}\ \ CC BY-NC-SA 4.0 · fentmates.org\qquad\small\thepage}}
% la resta: el símbol a l'esquerra i el curs a la dreta, amb el títol al centre com abans
\fancyhead[L]{\raisebox{-0.3\height}{\includegraphics[height=0.62cm]{\fm@marca}}}
\fancyhead[R]{\small\bfseries\color{cursfosc}\fm@cursnom}
\renewcommand{\headrule}{\color{curs!45}\hrule height \headrulewidth}
\makeatother

%   \fmcurs{1eso}
% i, al començament del document:
%   \capcalera{Matemàtiques}{1r ESO}{Fitxa d'exercicis · Unitat 2}{Divisibilitat}
% Opcional abans de \input: \def\fmneed{14} (espai mínim abans de cada secció).
% =====================================================================
\makeatletter
\def\fm@logodir{../../logo/pdf}
\providecommand{\fmneed}{10}
\def\fm@cursnom{}

% clau del curs, color i color del text damunt del color
\newcommand{\fm@curs}[3]{\@namedef{fm@color@#1}{#2}\@namedef{fm@text@#1}{#3}}
\fm@curs{1eso}{E6B84A}{navy}
\fm@curs{2eso}{3B82F6}{white}
\fm@curs{3eso}{FF7A50}{white}
\fm@curs{4eso-a}{A78BFA}{white}
\fm@curs{4eso-b}{6D28D9}{white}
\fm@curs{1batx-mi}{9B1C1C}{white}
\fm@curs{1batx-ccss}{EF6B6B}{white}
\fm@curs{2batx-mii}{1E6B47}{white}
\fm@curs{2batx-ccss}{2BB3A3}{white}

\newcommand{\fmcurs}[1]{%
  \@ifundefined{fm@color@#1}{\PackageError{fentmates}{Curs desconegut: #1}{}}{}%
  \edef\fm@tmp{\noexpand\definecolor{curs}{HTML}{\@nameuse{fm@color@#1}}%
               \noexpand\colorlet{curstext}{\@nameuse{fm@text@#1}}}\fm@tmp
  \colorlet{cursfosc}{curs!72!black}%
  \def\fm@logo{\fm@logodir/vertical-#1.pdf}%
  \def\fm@marca{\fm@logodir/marca-#1.pdf}}

% ---------- capçalera de la primera pàgina ----------
% #1 matèria (petit), #2 curs, #3 tipus de document i unitat, #4 títol
\newcommand{\capcalera}[4]{%
  \gdef\fm@cursnom{#2}%
  \par\noindent
  \begin{tcolorbox}[colback=curs,colframe=curs,boxrule=0pt,arc=2.2mm,width=3.5cm,height=2.15cm,
      valign=center,left=3mm,right=2mm,top=1mm,bottom=1mm,nobeforeafter,box align=center]
    \color{curstext}{\footnotesize #1}\par\vspace{1pt}{\fontsize{19}{21}\selectfont\bfseries #2\par}%
  \end{tcolorbox}\hspace{5mm}%
  \begin{minipage}[c]{\dimexpr\linewidth-3.5cm-5mm-2.9cm\relax}
    {\color{cursfosc}\small\bfseries\MakeUppercase{#3}}\par\vspace{2pt}
    {\color{navy}\fontsize{20}{23}\selectfont\bfseries #4\par}
  \end{minipage}\hfill
  \begin{minipage}[c]{2.6cm}\raggedleft\includegraphics[width=2.6cm]{\fm@logo}\end{minipage}\par
  \vspace{2pt}}

% ---------- barres de secció del color del curs ----------
\renewcommand{\seccio}[2]{\par\needspace{\fmneed\baselineskip}\bigskip
  \begin{tcolorbox}[colback=curs,colframe=curs,coltext=curstext,arc=2mm,boxrule=0pt,left=3mm,right=3mm,top=1.2mm,bottom=1.2mm]
  {\large\bfseries #1\quad #2}
  \end{tcolorbox}\vspace{-2pt}}

% ---------- estrelles de dificultat del color del curs (a Batxillerat no se'n posen) ----------
\@ifundefined{estrelles}{}{%
  \renewcommand{\estrelles}[1]{{\color{curs}\ifcase#1\or\faStar\faStar[regular]\faStar[regular]\or\faStar\faStar\faStar[regular]\or\faStar\faStar\faStar\fi}}}

% ---------- capçalera i peu de les pàgines ----------
% primera pàgina: logos institucionals al peu
\fancypagestyle{primera}{\fancyhf{}\renewcommand{\headrulewidth}{0pt}
  \fancyfoot[L]{\raisebox{-0.32\height}{\includegraphics[height=0.78cm]{logo-gva.png}}\hspace{5mm}%
    \raisebox{-0.32\height}{\includegraphics[height=0.78cm]{logo-vermellar.png}}}
  \fancyfoot[R]{\scriptsize\color{muted}\href{https://creativecommons.org/licenses/by-nc-sa/4.0/deed.ca}{\ccbyncsa}\ \ CC BY-NC-SA 4.0 · fentmates.org\qquad\small\thepage}}
% la resta: el símbol a l'esquerra i el curs a la dreta, amb el títol al centre com abans
\fancyhead[L]{\raisebox{-0.3\height}{\includegraphics[height=0.62cm]{\fm@marca}}}
\fancyhead[R]{\small\bfseries\color{cursfosc}\fm@cursnom}
\renewcommand{\headrule}{\color{curs!45}\hrule height \headrulewidth}
\makeatother

\fmcurs{2batx-mii}

\begin{document}
\thispagestyle{primera}

\capcalera{Matemàtiques II}{2n Batx}{Fitxa d'exercicis · Unitat 3}{Vectors i aplicacions. Rectes, plans i feixos de plans}
\vspace{10pt}
Nom i cognoms: \rule{8.4cm}{0.4pt}\hfill Grup: \rule{1.6cm}{0.4pt}\hfill Data: \rule{2.2cm}{0.4pt}
\vspace{8pt}

\begin{tcolorbox}[colback=soft,colframe=soft,arc=2mm,left=3mm,right=3mm,top=1.5mm,bottom=1.5mm]
\small Resol els exercicis al quadern justificant cada pas. Fes un esbós quan t'ajude a veure la situació i comprova sempre els resultats (que el punt trobat és dels plans, que el vector normal és perpendicular als vectors directors\ldots). Els exercicis marcats amb {\color{teal}\bfseries\faUniversity\ PAU} són de les proves d'accés a la universitat d'altres comunitats autònomes (2024--2026), traduïts al valencià; n'hem tret els apartats d'angles i distàncies, que són de la unitat següent.
\end{tcolorbox}

% =====================================================
\seccio{1}{Vectors a l'espai i aplicacions}

\ex{1} Siguen $A(1,\,-2,\,3)$, $B(3,\,-1,\,1)$ i $C(-1,\,2,\,7)$.\begin{tasks}(2)
  \task Calcula $\overrightarrow{AB}$, $\overrightarrow{AC}$ i els seus mòduls.
  \task Troba un vector unitari amb la direcció i el sentit de $\overrightarrow{AC}$.
  \task Troba un vector de mòdul 9 paral·lel a $\overrightarrow{AB}$ i de sentit contrari.
  \task Calcula $3\overrightarrow{AB}+\overrightarrow{AC}$.
\end{tasks}

\ex{1} El punt mig del segment $AB$ és $M(1,\,-2,\,4)$ i $A(3,\,0,\,-1)$. Calcula $B$. Troba també el simètric de $P(2,\,-3,\,1)$ respecte del punt $Q(0,\,1,\,-2)$.\recorda{El punt mig de $A$ i $B$ és $M=\frac{A+B}{2}$; per tant, $B=2M-A$. El simètric d'un punt respecte d'un altre es calcula igual.}

\ex{2} Divideix en tres parts iguals el segment d'extrems $A(-2,\,1,\,5)$ i $B(4,\,-5,\,-1)$.

\ex{2} Troba $a$ i $b$ perquè els punts $A(1,\,0,\,2)$, $B(2,\,1,\,0)$ i $C(a,\,3,\,b)$ estiguen alineats.

\ex{2} Els punts $A(1,\,2,\,-1)$, $B(3,\,0,\,2)$ i $C(2,\,-1,\,4)$ són tres vèrtexs consecutius del paral·lelogram $ABCD$. Calcula el quart vèrtex $D$ i el centre del paral·lelogram.

\ex{2} El baricentre del triangle $ABC$ és $G(1,\,2,\,0)$, amb $A(2,\,1,\,-3)$ i $B(-1,\,4,\,2)$. Troba el vèrtex $C$.

\ex{2} Troba el punt $P$ del segment $AB$, amb $A(1,\,-1,\,2)$ i $B(7,\,5,\,-4)$, que compleix $\overline{AP}=2\,\overline{PB}$.

\ex{2} \pau{Castella-la Manxa, extraordinària 2026} Es consideren els punts $A(1,\,2,\,-1)$, $B(3,\,6,\,1)$ i $C(m,\,m,\,m-2)$, on $m$ és un nombre real. Determina el valor de $m$ perquè els punts $A$, $B$ i $C$ estiguen alineats.

\ex{2} \pau{Catalunya, extraordinària 2026} L'Aina vol fer un croquis en 3D d'una capsa amb un programa informàtic. Situa un costat de la tapa des del vèrtex $A(1,\,1,\,1)$ fins al vèrtex $B(4,\,1,\,5)$, i l'altre costat de la tapa des del vèrtex $A$ fins al vèrtex $C(1,\,10,\,1)$. Quines són les coordenades del quart vèrtex de la tapa? Calcula les longituds dels dos costats de la tapa.


% =====================================================
\seccio{2}{Dependència lineal i bases}

\ex{1} Estudia si els vectors $\vec u=(1,\,2,\,-1)$, $\vec v=(2,\,1,\,3)$ i $\vec w=(0,\,3,\,-5)$ són linealment independents. Si no ho són, escriu $\vec w$ com a combinació lineal de $\vec u$ i $\vec v$.

\ex{2} Considera els vectors $\vec u=(1,\,1,\,1)$, $\vec v=(1,\,a,\,1)$ i $\vec w=(1,1,a^2)$.\begin{tasks}(1)
  \task Per a quins valors de $a$ formen una base de $\mathbb R^3$?
  \task Calcula el rang del conjunt $\{\vec u,\vec v,\vec w\}$ per als altres valors de $a$.
\end{tasks}

\ex{2} Comprova que $\mathcal B=\{(1,\,1,\,0),(0,\,1,\,1),(1,\,0,\,1)\}$ és una base i calcula les coordenades del vector $\vec x=(3,\,1,\,2)$ en aquesta base.

\ex{2} Es pot escriure $(1,\,2,\,3)$ com a combinació lineal de $(1,\,1,\,0)$ i $(0,\,1,\,1)$? I $(2,\,5,\,3)$?

\ex{3} Siga $\{\vec u,\vec v,\vec w\}$ una base. Estudia si són linealment independents:\begin{tasks}(2)
  \task $\{\vec u+\vec v,\ \vec v+\vec w,\ \vec u+\vec w\}$
  \task $\{\vec u-\vec v,\ \vec v-\vec w,\ \vec w-\vec u\}$
\end{tasks}\recorda{Escriu cada vector nou amb coordenades en la base $\{\vec u,\vec v,\vec w\}$ i calcula el determinant d'aquestes coordenades.}


% =====================================================
\seccio{3}{Producte escalar}

\ex{1} Donats $\vec u=(1,\,-2,\,2)$ i $\vec v=(3,\,0,\,-4)$, calcula:\begin{tasks}(2)
  \task $\vec u\cdot\vec v$, $|\vec u|$ i $|\vec v|$.
  \task L'angle que formen.
  \task La projecció escalar de $\vec u$ sobre $\vec v$.
  \task El vector projecció de $\vec u$ sobre $\vec v$.
\end{tasks}

\ex{2} Troba els vectors de mòdul $\sqrt{29}$ que són perpendiculars alhora a $(1,\,2,\,-1)$ i a $(2,\,0,\,1)$ (sense fer servir el producte vectorial).

\ex{2} Troba $m$ perquè els vectors $(2,\,m,\,1)$ i $(1,\,1,\,2)$ formen un angle de $60^\circ$.

\ex{2} Calcula els angles que forma el vector $\vec v=(2,\,-1,\,2)$ amb els eixos de coordenades. Pot haver-hi un vector que forme $45^\circ$ amb l'eix $OX$ i $60^\circ$ amb l'eix $OY$? Quin angle formaria amb l'eix $OZ$?

\ex{3} Sabem que $|\vec u|=3$, $|\vec v|=4$ i $\vec u\cdot\vec v=6$. Calcula:\begin{tasks}(2)
  \task L'angle que formen $\vec u$ i $\vec v$.
  \task $|\vec u+\vec v|$ i $|\vec u-\vec v|$.
  \task $(2\vec u-\vec v)\cdot(\vec u+3\vec v)$. Què hi observes?
\end{tasks}\recorda{$|\vec a|^2=\vec a\cdot\vec a$, i el producte escalar és distributiu.}

\ex{2} Demostra que el triangle de vèrtexs $A(1,\,1,\,1)$, $B(3,\,2,\,-1)$ i $C(2,\,3,\,3)$ és rectangle i isòsceles, i calcula'n els angles.

\ex{1} Troba $m$ perquè els vectors $(m,\,1,\,3)$ i $(2,\,m,\,-1)$ siguen perpendiculars.


% =====================================================
\seccio{4}{Producte vectorial i producte mixt. Àrees i volums}

\ex{1} Siguen $\vec u=(2,\,-1,\,3)$ i $\vec v=(1,\,0,\,-2)$. Calcula $\vec u\times\vec v$ i $\vec v\times\vec u$, i comprova que $\vec u\times\vec v$ és perpendicular a $\vec u$ i a $\vec v$.

\ex{2} Donats $A(1,\,0,\,1)$, $B(2,\,1,\,3)$ i $C(0,\,2,\,1)$, troba els vectors unitaris perpendiculars al pla que determinen i l'àrea del triangle $ABC$.

\ex{2} Els punts $A(1,\,1,\,0)$, $B(3,\,2,\,1)$ i $D(0,\,3,\,2)$ són vèrtexs del paral·lelogram $ABCD$. Calcula el vèrtex $C$ i l'àrea del paral·lelogram.

\ex{2} Troba $k$ perquè el triangle de vèrtexs $A(1,\,0,\,0)$, $B(0,\,1,\,0)$ i $C(0,\,0,\,k)$ tinga àrea $\tfrac32$.

\ex{2} Troba $a$ perquè el tetraedre de vèrtexs $A(1,\,1,\,1)$, $B(3,\,1,\,0)$, $C(1,\,4,\,2)$ i $D(1,\,a,\,3)$ tinga volum 2.

\ex{2} Troba $a$ perquè els punts $A(1,\,2,\,1)$, $B(2,\,3,\,0)$, $C(1,\,0,\,3)$ i $D(a,\,1,\,a)$ siguen coplanaris.

\ex{3} Sabem que $[\vec u,\vec v,\vec w]=5$. Calcula:\begin{tasks}(3)
  \task $[\vec w,\vec v,\vec u]$
  \task $[2\vec u,\ \vec v-\vec w,\ 3\vec w]$
  \task $[\vec u+\vec v,\ \vec v,\ \vec w]$
\end{tasks}\recorda{El producte mixt és un determinant: aplica'n les propietats (permutar files, treure factors, sumar a una fila una combinació de les altres).}

\ex{2} \pau{Andalusia, ordinària 2026} Considera els punts $A(1,\,2,\,0)$, $B(2,\,m,\,1)$, $C(3,\,4,\,2)$ i $D(1,\,-1,\,m)$.\begin{tasks}(1)
  \task Troba els valors de $m$ per als quals els punts són coplanaris.
  \task Per a $m=1$, calcula l'àrea del triangle de vèrtexs $A$, $B$, $C$ i el volum del tetraedre de vèrtexs $A$, $B$, $C$, $D$.
\end{tasks}

\ex{2} \pau{Aragó, ordinària 2024} Donats els punts $P_1(-2,\,1,\,1)$, $P_2(0,\,a,\,-2)$, $P_3(-1,\,1,\,-1)$ i $P_4(1,\,3,\,-3)$:\begin{tasks}(1)
  \task Calcula els valors de $a$ perquè el tetraedre de vèrtexs $P_1$, $P_2$, $P_3$ i $P_4$ tinga volum $\tfrac13$.
  \task Calcula el valor de $a$ perquè els quatre punts siguen coplanaris.
\end{tasks}

\ex{2} \pau{Astúries, ordinària 2025} Es consideren els punts $A(1,\,2,\,3)$, $B(-2,\,1,\,4)$, $C(3,\,0,\,5)$ i $D(0,\,-1,\,2)$.\begin{tasks}(1)
  \task Estudia si els quatre punts són d'un mateix pla.
  \task Calcula l'àrea del triangle de vèrtexs $A$, $B$ i $C$.
  \task Calcula el volum del tetraedre que formen els quatre punts.
\end{tasks}

\ex{2} \pau{Castella-la Manxa, ordinària 2026} Donats els vectors $\vec u=(1,\,-1,\,a)$, $\vec v=(2,\,a,\,1)$ i $\vec w=(-1,\,1,\,-1)$:\begin{tasks}(1)
  \task Troba els valors reals de $a$ perquè els tres vectors siguen coplanaris. Justifica la resposta.
  \task Troba els valors reals de $a$ perquè el volum del paral·lelepípede format pels tres vectors siga $10\ u^3$.
\end{tasks}

\ex{2} \pau{Madrid, extraordinària 2025} Siguen els punts $A(1,\,1,\,2)$, $B(2,\,-1,\,0)$, $C(-2,\,0,\,3)$ i $D(2,\,-3,\,-1)$.\begin{tasks}(1)
  \task Comprova que els punts no són coplanaris i calcula el volum del tetraedre que determinen.
  \task Calcula l'àrea de la cara $ABC$ del tetraedre i la longitud de l'altura del tetraedre que parteix del vèrtex $D$.
\end{tasks}\recorda{El volum d'una piràmide és un terç de l'àrea de la base per l'altura.}

\ex{2} \pau{Navarra, ordinària 2024} Els punts $A(4,\,-2,\,-3)$, $B(2,\,-1,\,1)$ i $C(0,\,-3,\,-1)$ són vèrtexs d'un rombe.\begin{tasks}(1)
  \task Troba el quart vèrtex del rombe.
  \task Calcula l'àrea del rombe.
\end{tasks}

\ex{2} \pau{Aragó, extraordinària 2024} Siguen $A(1,\,2,\,3)$, $B(1,\,0,\,-1)$ i $C(2,\,2,\,2)$ tres punts de l'espai, $\vec v_1$ el vector que va de $A$ a $B$, $\vec v_2$ el vector que va de $B$ a $C$ i $\vec v_3$ el vector que va de $C$ a $A$.\begin{tasks}(1)
  \task Estudia si els vectors $\vec v_1$, $\vec v_2$ i $\vec v_3$ són linealment independents.
  \task Calcula l'àrea del triangle de vèrtexs $A$, $B$ i $C$.
\end{tasks}


% =====================================================
\seccio{5}{Equacions de rectes i plans}

\ex{1} Escriu les equacions vectorial, paramètriques, contínua i implícites de la recta $r$ que passa per $A(2,\,-1,\,3)$ i $B(4,\,0,\,1)$. Els punts $C(0,\,-2,\,5)$ i $D(6,\,1,\,0)$ són de $r$?

\ex{1} Escriu en forma paramètrica la recta $r:\left\{\begin{aligned}x-y+z&=2\\2x+y-z&=1\end{aligned}\right.$

\ex{1} Troba les equacions paramètriques i general del pla que passa per $A(1,\,0,\,2)$, $B(2,\,1,\,0)$ i $C(0,\,3,\,1)$.

\ex{2} Troba l'equació general del pla que passa per $P(2,\,-1,\,1)$ i és paral·lel a les rectes \mbox{$r:\dfrac{x-1}{2}=\dfrac y1=\dfrac{z+1}{-1}$} i $s:(x,\,y,\,z)=(0,\,1,\,0)+\mu(1,\,0,\,3)$.

\ex{2} Comprova que les rectes \mbox{$r:\dfrac{x-1}{1}=\dfrac{y+1}{2}=\dfrac{z}{-1}$} i \mbox{$s:\dfrac{x+1}{-2}=\dfrac{y-2}{-4}=\dfrac{z-1}{2}$} són paral·leles i troba el pla que les conté.

\ex{2} Escriu l'equació segmentària i la general del pla que talla els eixos en $(2,\,0,\,0)$, $(0,\,-3,\,0)$ i $(0,\,0,\,6)$. Quin volum té el tetraedre que formen el pla $4x-2y+3z=12$ i els plans coordenats?

\ex{2} Troba el pla que passa per l'origen i és perpendicular a la recta $r:\left\{\begin{aligned}x+y&=1\\y-z&=2\end{aligned}\right.$, i la recta que passa per $A(1,\,2,\,3)$ i és perpendicular al pla $2x-y+3z=5$.

\ex{1} Passa a forma general el pla $(x,\,y,\,z)=(1,\,0,\,-1)+\lambda(1,\,2,\,0)+\mu(0,\,1,\,3)$ i troba $a$, $b$ perquè $P(a,\,1,\,b)$ siga de la recta \mbox{$\dfrac{x-2}{3}=\dfrac{y+1}{2}=\dfrac{z-1}{-1}$}.

\ex{1} \pau{Andalusia, ordinària 2025} Considera el punt $P(1,\,1,\,1)$ i la recta \mbox{$r:\dfrac{x-1}{1}=\dfrac{y-2}{2}=\dfrac{z-3}{2}$}. Troba el pla que passa pel punt $P$ i conté la recta $r$.

\ex{2} \pau{Aragó, ordinària 2025}\begin{tasks}(1)
  \task Troba les equacions paramètriques de la recta $s$ que passa pel punt $P(4,\,-3,\,0)$ i és perpendicular al pla $\pi:x-2y+z-1=0$.
  \task Troba l'equació del pla que conté el punt $Q(1,\,2,\,3)$ i la recta $r:\left\{\begin{aligned}x+y+z&=1\\x-2y&=0\end{aligned}\right.$
\end{tasks}

\ex{2} \pau{Cantàbria, extraordinària 2024} Siguen $A(0,\,3,\,2)$, $B(4,\,1,\,3)$, $C(2,\,3,\,4)$ i $D(0,\,1,\,2)$ els vèrtexs d'un tetraedre.\begin{tasks}(1)
  \task Obtín l'equació vectorial del pla determinat pels punts $A$, $B$ i $C$.
  \task Calcula el volum del tetraedre.
\end{tasks}

\ex{2} \pau{Castella-la Manxa, extraordinària 2026} Es consideren els punts $A(1,\,1,\,0)$, $B(2,\,0,\,1)$, $C(0,\,2,\,1)$ i $D(a,\,1,\,2)$, amb $a$ un nombre real.\begin{tasks}(1)
  \task Determina l'equació del pla que passa pels punts $A$, $B$ i $C$.
  \task Determina els valors de $a$ perquè el tetraedre format pels quatre punts tinga volum 3 unitats cúbiques.
\end{tasks}

\ex{2} \pau{Castella i Lleó, extraordinària 2026} Els punts $A(1,\,1,\,1)$ i $B(0,\,2,\,0)$ són vèrtexs consecutius del paral·lelogram $ABCD$, que té el centre en $E(0,\,0,\,1)$.\begin{tasks}(1)
  \task Calcula les coordenades dels altres vèrtexs.
  \task Troba l'equació general del pla que conté el paral·lelogram.
  \task Troba el volum del tetraedre de vèrtexs $A$, $B$, $E$ i l'origen de coordenades $O$.
\end{tasks}

\ex{2} \pau{Catalunya, ordinària 2026} Considera els punts $P(1,\,0,\,-1)$, $Q(3,\,-2,\,0)$ i $R(1,\,1,\,1)$.\begin{tasks}(1)
  \task Calcula l'equació del pla que conté els punts $P$, $Q$ i $R$.
  \task Comprova que l'àrea del triangle $PQR$ és $\tfrac{3\sqrt5}{2}\ u^2$.
  \task Determina les condicions que han de complir les coordenades d'un quart punt $S(x,\,y,\,z)$ perquè $P$, $Q$, $R$ i $S$ formen un tetraedre de volum 1.
\end{tasks}

\ex{2} \pau{Extremadura, extraordinària 2026} Donats els punts $A(4,\,3,\,3)$, $B(3,\,3,\,4)$, $C(1,\,0,\,4)$ i $D(4,\,0,\,1)$:\begin{tasks}(1)
  \task Justifica de forma raonada que els punts $A$, $B$, $C$ i $D$ estan en el mateix pla i troba l'equació del pla que determinen.
  \task Calcula el valor de $k$ perquè el punt $P(7,\,k,\,2)$ siga de la recta $r$ que passa pels punts $A$ i $C$.
\end{tasks}

\ex{2} \pau{Extremadura, ordinària 2026}\begin{tasks}(1)
  \task Donats els vectors $\vec u=(0,\,3,\,-4)$ i $\vec v=(12,\,4,\,3)$, comprova que el paral·lelogram que formen és un rectangle i calcula'n l'àrea.
  \task Troba l'equació del pla format pels vectors $\vec u$ i $\vec v$ i el punt $(1,\,0,\,0)$.
  \task Obtín un vector perpendicular als vectors $\vec u$ i $\vec v$ de mòdul 13.
\end{tasks}

\ex{2} \pau{Madrid, ordinària 2025} Donats la recta \mbox{$r:\dfrac{x-1}{2}=\dfrac y0=\dfrac{z-2}{1}$} i el pla $\pi:x+2y-3z=1$, troba una equació del pla que conté $r$ i és perpendicular a $\pi$.


% =====================================================
\seccio{6}{Posicions relatives de rectes i plans}

\ex{1} Estudia la posició relativa de \mbox{$r:\dfrac{x-1}{2}=\dfrac{y+2}{-1}=\dfrac z3$} i $s:(x,\,y,\,z)=(3,\,-3,\,4)+\mu(4,\,-2,\,6)$. Si són coplanàries, troba el pla que les conté.

\ex{1} Estudia la posició relativa de $r:(x,\,y,\,z)=(1,\,0,\,2)+\lambda(2,\,1,\,-1)$ i $s:(x,\,y,\,z)=(0,\,3,\,1)+\mu(1,\,-2,\,1)$.

\ex{2} Troba $a$ perquè les rectes \mbox{$r:\dfrac{x-1}{1}=\dfrac{y-2}{2}=\dfrac{z-a}{1}$} i $s:(x,\,y,\,z)=(0,\,1,\,3)+\mu(1,\,0,\,1)$ es tallen, i calcula el punt de tall.

\ex{2} Estudia, segons els valors de $k$, la posició relativa de $r:(x,\,y,\,z)=(2,\,-1,\,0)+\lambda(1,\,1,\,2)$ i $\pi:3x-y-z+k=0$.

\ex{2} Estudia, segons els valors de $m$, la posició de la recta $r:\left\{\begin{aligned}x+y-z&=1\\2x-y+z&=2\end{aligned}\right.$ i el pla $\pi:x+my+z=2$. Troba el punt de tall quan $m=1$.

\ex{2} Estudia la posició relativa dels plans $\pi_1:2x-y+3z=1$ i $\pi_2:4x+ay+6z=b$ segons els valors de $a$ i $b$. Per a $a=1$, $b=2$, troba la recta en què es tallen.

\ex{3} Troba la recta que passa per $P(1,\,1,\,1)$, és paral·lela al pla $\pi:x+y+z=0$ i talla la recta $s:(x,\,y,\,z)=(2,\,1,\,0)+\mu(0,\,1,\,1)$.\recorda{La recta buscada està continguda en el pla paral·lel a $\pi$ que passa per $P$. Talla $s$ en el punt on $s$ talla aquest pla.}

\ex{2} \pau{Astúries, extraordinària 2024} Es consideren els punts $A(1,\,1,\,1)$, $B(1,\,0,\,2)$, $C(-1,\,1,\,3)$ i $D(-1,\,0,\,1)$.\begin{tasks}(1)
  \task Estudia si hi ha un pla que continga els quatre punts.
  \task Calcula la recta $r$ que passa per $D$ i és perpendicular al pla $\pi$ que conté $A$, $B$ i $C$.
  \task Calcula el punt $P$ d'intersecció de $r$ i $\pi$.
\end{tasks}

\ex{2} \pau{Canàries, extraordinària 2025} Es tenen les rectes $r:\left\{\begin{aligned}x+y-2z&=3\\2x+y-6z&=2\end{aligned}\right.$ \ i \ \mbox{$s:x-1=\dfrac y3=\dfrac{z+1}{2}$}.\begin{tasks}(1)
  \task Comprova que $r$ i $s$ són coplanàries.
  \task Troba l'equació del pla que les conté.
\end{tasks}

\ex{2} \pau{Canàries, ordinària 2026} Es coneixen el pla $\pi:2x+3y=0$ i el punt $P(1,\,-5,\,1)$.\begin{tasks}(1)
  \task Troba la recta $r$ perpendicular al pla $\pi$ que passa pel punt $P$.
  \task Troba la posició relativa de la recta $s:\left\{\begin{aligned}-x+3y-3z&=-3\\4x+9y-2z&=-2\end{aligned}\right.$ \ amb el pla $\pi$.
\end{tasks}

\ex{3} \pau{Canàries, extraordinària 2024} Considera les rectes $r_1:\left\{\begin{aligned}x-3y+2z+2&=0\\2x+y-3z&=3\end{aligned}\right.$ \ i \ \mbox{$r_2:\dfrac{1-x}{2}=y=\dfrac{1-z}{2}$}.\begin{tasks}(1)
  \task Estudia la posició relativa de les dues rectes.
  \task Troba l'equació de la recta $s$ que té direcció perpendicular a les dues rectes i passa per $P\left(0,\,\tfrac12,\,0\right)$. Calcula el punt de tall de $s$ amb $r_1$.
\end{tasks}

\ex{2} \pau{Cantàbria, ordinària 2024} Considera la recta $r:\left\{\begin{aligned}x+y+z+5&=0\\x+2y-z&=0\end{aligned}\right.$ \ i el pla $\pi:2x+y-az=3$. Raona si és possible donar algun valor al paràmetre $a$ perquè:\begin{tasks}(1)
  \task la recta estiga continguda en el pla;
  \task la recta i el pla siguen paral·lels;
  \task la recta i el pla es tallen. En cas afirmatiu, dona un valor de $a$ i calcula on es tallen.
\end{tasks}

\ex{2} \pau{Castella i Lleó, ordinària 2025} Siguen les rectes \mbox{$r:\dfrac x4=\dfrac{y-1}{-2}=\dfrac z2$} i \mbox{$s:x-1=\dfrac{y-m}{m-1}=\dfrac{z-3}{3}$}.\begin{tasks}(1)
  \task Comprova que les rectes $r$ i $s$ es tallen per a qualsevol valor de $m$.
  \task Per a $m=6$, troba el punt d'intersecció de les rectes $r$ i $s$.
\end{tasks}

\ex{2} \pau{Castella i Lleó, extraordinària 2025} Siguen el pla $\pi:5x+my+z=2$ i la recta $r:(x,\,y,\,z)=(1,\,1,\,0)+t(-1,\,-1,\,2)$.\begin{tasks}(1)
  \task Determina la posició relativa de $r$ i $\pi$ en funció de $m$.
  \task Per a $m=1$, calcula el pla $\pi'$ que conté $r$ i és perpendicular a $\pi$.
\end{tasks}

\ex{1} \pau{Galícia, extraordinària 2025} Donades les rectes\[r:\dfrac{x-1}{2}=\dfrac{y-2}{-1}=\dfrac{z-1}{1}\qquad\text{i}\qquad s:\dfrac{x-2}{4}=\dfrac{y-1}{-2}=\dfrac{z-1}{2}\]\begin{tasks}(1)
  \task Calcula la posició relativa de les rectes $r$ i $s$.
  \task Obtín l'equació del pla que conté les rectes $r$ i $s$.
\end{tasks}

\ex{2} \pau{Galícia, extraordinària 2024} Considera el pla $\pi:x+2y-2z=0$ i la recta $r$ que passa pels punts $A(2,\,1,\,2)$ i $B(0,\,1,\,1)$.\begin{tasks}(1)
  \task Estudia la posició relativa de la recta $r$ i el pla $\pi$.
  \task Obtín l'equació implícita o general del pla que conté $r$ i és perpendicular a $\pi$.
\end{tasks}

\ex{1} \pau{Galícia, ordinària 2024} Considera el pla $\pi:4x+2y+bz=2$ i la recta \mbox{$r:\dfrac{x-2}{3}=\dfrac{y-c}{2}=\dfrac{z-3}{4}$}, on $b$ i $c$ són paràmetres reals. Calcula els valors de $b$ i $c$ perquè la recta $r$ estiga continguda en $\pi$.

\ex{2} \pau{Múrcia, extraordinària 2025} Considera la recta $r:\left\{\begin{aligned}x+y+z&=1\\-x-2y+z&=0\end{aligned}\right.$ \ i el pla $\pi:2x+y+az=b$.\begin{tasks}(1)
  \task Determina els valors de $a$ i $b$ perquè el pla contínga la recta.
  \task Per a quins valors de $a$ i $b$ la recta talla el pla? Troba el punt de tall en el cas $a=0$, $b=7$.
\end{tasks}

\ex{2} \pau{Navarra, extraordinària 2024} Es consideren el pla $\pi:2x+y-z-5=0$, la recta $r:\left\{\begin{aligned}x+2z+3&=0\\-x-y+z+4&=0\end{aligned}\right.$ \ i els punts $A(3,\,2,\,-1)$ i $B(1,\,1,\,-1)$. Siga $C$ la intersecció de la recta i el pla.\begin{tasks}(1)
  \task Demostra que els punts $A$, $B$ i $C$ no estan alineats.
  \task Calcula l'àrea del triangle que formen els tres punts.
\end{tasks}

\ex{3} \pau{Navarra, ordinària 2025} Calcula l'equació contínua de la recta $t$ que passa pel punt $P(2,\,0,\,-1)$ i talla les rectes $s:\left\{\begin{aligned}2x+y-3z-6&=0\\2x-3z-8&=0\end{aligned}\right.$ \ i \ \mbox{$r:\dfrac{x+1}{2}=\dfrac y1=\dfrac{z+2}{1}$}.\recorda{La recta $t$ és la intersecció del pla que conté $P$ i $s$ amb el pla que conté $P$ i $r$.}

\ex{2} \pau{País Basc, extraordinària 2024} Es consideren la recta $r:\left\{\begin{aligned}2x-y+z&=0\\x-y+4z&=1\end{aligned}\right.$ \ i el pla $\pi:2x-3y+Az=10$.\begin{tasks}(1)
  \task Calcula el valor del paràmetre $A$ perquè la recta $r$ i el pla $\pi$ siguen paral·lels.
  \task Si $A=21$, calcula la intersecció del pla $\pi$ i la recta $r$.
\end{tasks}

\ex{1} \pau{Aragó, ordinària 2026} Siguen els plans $\pi:x-2y+4z+1=0$ i $\pi_d:2x-3y+6z=d$, amb $d\in\mathbb R$.\begin{tasks}(1)
  \task Hi ha algun valor de $d$ per al qual els plans $\pi$ i $\pi_d$ no es tallen?
  \task Calcula, per a $d=-2$, els punts de tall entre els plans $\pi$ i $\pi_d$.
\end{tasks}


% =====================================================
\seccio{7}{Posició relativa de tres plans}

\ex{2} Estudia, segons el valor de $a$, la posició relativa dels plans $\pi_1:x+y-z=1$, $\pi_2:2x-y+z=2$ i $\pi_3:x-2y+2z=a$. Quan tinguen una recta comuna, troba-la.

\ex{3} Discuteix, segons els valors de $k$, la posició relativa dels plans $kx+y+z=k$, \ $x+ky+z=k$, \ $x+y+kz=k$, i troba el punt comú quan n'hi haja un de sol.

\ex{2} Troba $a$ i $b$ perquè els plans $x+y+z=2$, $x-y+2z=1$ i $3x+y+az=b$ tinguen una recta en comú. Què passa si $a=4$ i $b\neq5$? I si $a\neq4$?

\ex{2} Donats $\pi_1:x-y=0$ i $\pi_2:y-z=0$, escriu un pla $\pi_3$ tal que els tres plans\begin{tasks}(2)
  \task formen un prisma;
  \task es tallen en un sol punt.
\end{tasks}

\ex{3} \pau{La Rioja, extraordinària 2026} Determina els valors de $a$ perquè els plans $\pi_1:x+y+z=1$, \ $\pi_2:(a-1)x+y+z=a$ \ i \ $\pi_3:x+(a-1)y-z=0$\begin{tasks}(1)
  \task es tallen en un punt. En aquest cas, determina el punt de tall.
  \task es tallen en una recta. En aquest cas, determina la recta en forma paramètrica.
\end{tasks}


% =====================================================
\seccio{8}{Feixos de plans}

\ex{1} Siga $r:\left\{\begin{aligned}x-y+z&=1\\2x+y&=3\end{aligned}\right.$. Escriu el feix de plans d'eix $r$ i troba el pla del feix que passa per l'origen.

\ex{2} Amb la recta $r$ de l'exercici anterior, troba el pla que conté $r$ i és paral·lel a l'eix $OZ$, i el pla que conté $r$ i és perpendicular a $\pi:x+2y-z=0$.

\ex{2} Troba el pla paral·lel a $2x-y+2z=1$ que passa per $A(1,\,-1,\,2)$, i els plans d'aquest feix que formen amb els plans coordenats un tetraedre de volum 9.

\ex{2} Escriu la recta \mbox{$r:\dfrac{x-1}{2}=\dfrac y1=\dfrac{z+1}{3}$} com a intersecció de dos plans i troba, amb el feix de plans, el pla que conté $r$ i passa per $P(1,\,1,\,1)$.

\ex{2} \pau{Múrcia, extraordinària 2026} Determina l'equació del pla que conté el punt $P(-2,\,2,\,-2)$ i la recta $r$ donada com a intersecció de dos plans: \ $r:\left\{\begin{aligned}2x+y-z&=4\\x+y-2z&=3\end{aligned}\right.$

\ex{2} \pau{País Basc, ordinària 2025} Es consideren les rectes $r_1:\left\{\begin{aligned}x+y-2z&=0\\2x-3y+z&=1\end{aligned}\right.$ \ i \ $r_2:\left\{\begin{aligned}x&=3t\\y&=1-2t\\z&=2+t\end{aligned}\right.$ \ Calcula l'equació del pla que conté $r_1$ i passa pel punt d'intersecció del pla $\pi:x-3y-2z+7=0$ i la recta $r_2$.


% =====================================================
\seccio{9}{Problemes i qüestions}

\ex{2} Raona si són certes o falses les afirmacions següents:\begin{tasks}(1)
  \task Si $\vec u\cdot\vec v=0$ i $\vec u\times\vec v=\vec 0$, aleshores algun dels dos vectors és nul.
  \task Tres vectors no nuls i perpendiculars dos a dos formen una base.
  \task Si $[\vec u,\vec v,\vec w]=0$, algun dels tres vectors és nul.
  \task Dues rectes que no tenen cap punt en comú són paral·leles.
  \task Si $\vec u\times\vec v=\vec u\times\vec w$ i $\vec u\neq\vec 0$, aleshores $\vec v=\vec w$.
\end{tasks}

\ex{2} Un cub té un vèrtex en l'origen, tres arestes sobre els semieixos positius i aresta 2. Troba l'equació del pla que passa pels vèrtexs $A(2,\,0,\,0)$, $B(0,\,2,\,0)$ i $C(0,\,0,\,2)$, l'àrea del triangle $ABC$ i el punt on la diagonal del cub que ix de l'origen talla aquest pla.

\ex{2} Un dron ix del punt $P(1,\,2,\,0)$ i vola en línia recta en la direcció del vector $(2,\,1,\,3)$. A la seua trajectòria hi ha la teulada inclinada d'una nau, continguda en el pla $x+y+z=12$, i un pal vertical que passa pel punt $(5,\,4,\,0)$. Amb quin dels dos obstacles topa primer?

\ex{2} \pau{Aragó, extraordinària 2026} Estem reparant una cadira metàl·lica que coixejava, perquè els extrems inferiors de les seues quatre potes no recolzaven alhora en un terra pla. La cadira està tombada sobre una taula de treball de 120 cm d'altura i té quatre potes de 45 cm de longitud cadascuna. Després de la reparació, els extrems inferiors de les potes estan situats en els punts $A(-2,\,3,\,120)$, $B(20,\,42,\,120)$, $C(-1,\,4,\,165)$ i $D(20,\,43,\,165)$. Quan baixem la cadira de la taula i la deixem en el terra, coixejarà o recolzarà les quatre potes sobre el terra pla del taller?

\ex{2} \pau{Illes Balears, extraordinària 2026} En un estudi d'astronomia es modelitza una constel·lació artificial de satèl·lits que orbiten al voltant de la Terra. Prenem un sistema de coordenades amb l'origen en el centre de la Terra (les unitats són desenes de milers de quilòmetres) i hi situem cinc satèl·lits: $A(1,\,2,\,3)$, $B(3,\,2,\,1)$, $C(5,\,4,\,3)$, $D(2,\,1,\,1)$ i $E(4,\,3,\,1)$. Diem que uns satèl·lits formen un pla orbital si, juntament amb el centre de la Terra, són coplanaris.\begin{tasks}(1)
  \task Determina si els satèl·lits $A$, $B$ i $C$ formen un pla orbital.
  \task Proporciona un punt $F$ (diferent de l'origen) on es podria col·locar un sisè satèl·lit perquè els satèl·lits $D$, $E$ i $F$ formen un altre pla orbital.
\end{tasks}

\ex{2} \pau{Illes Balears, ordinària 2026} En un estudi de gravació s'han penjat del sostre tres micròfons per a gravar un concert acústic. Prenem un sistema de referència amb l'origen en el terra de la sala; el sostre és a 7 m d'altura. El micròfon $M_1$ té coordenades $x=2$, $y=3$ i penja d'un cable de 1 m; el micròfon $M_2$ té $x=4$, $y=4$ i penja d'un cable de 2 m, i el micròfon $M_3$ té $x=5$, $y=1$ i penja d'un cable de 3 m. Determina les coordenades dels tres micròfons i calcula l'equació del pla que els conté.

\ex{2} \pau{La Rioja, ordinària 2026} En un xiringuito de platja s'instal·la un tendal per als clients. El tendal es col·loca en els punts de suport $\left(3,\,8,\,\tfrac{17}{2}\right)$, $\left(4,\,6,\,\tfrac{37}{4}\right)$, $(2,\,8,\,8)$ i $(s,\,16,\,8)$, amb $s\in\mathbb R$.\begin{tasks}(1)
  \task Determina l'equació del pla que conté el tendal.
  \task Calcula el valor de $s$ perquè els quatre punts siguen coplanaris.
\end{tasks}

\vfill
\begin{center}\color{muted}\small\faCheckCircle[regular]\ \ Ànim! Comprova els teus resultats amb les eines 3D dels apunts interactius.\end{center}

\end{document}
